
\documentclass[11pt]{article}

\usepackage{style}


\setlength{\droptitle}{-5em} %% Don't touch

% %%%%%%%%%%%%%%%%%%%%%%%%%%%%%%%%%%%%%%%%%%%%%%%%%%%%%%%%%%
% SET THE TITLE
% %%%%%%%%%%%%%%%%%%%%%%%%%%%%%%%%%%%%%%%%%%%%%%%%%%%%%%%%%%
% TITLE:
\title{Lasso Learning and Endogenous Return Predictability}

\author{
Tommaso Di Francesco\thanks{University of Bonn}
\and
Jonathan Seim\thanks{University of Cologne}
}
% DATE:
\date{\today}

% %%%%%%%%%%%%%%%%%%%%%%%%%%%%%%%%%%%%%%%%%%%%%%%%%%%%%%%%%%
% %%%%%%%%%%%%%%%%%%%%%%%%%%%%%%%%%%%%%%%%%%%%%%%%%%%%%%%%%%
\begin{document}
% %%%%%%%%%%%%%%%%%%%%%%%%%%%%%%%%%%%%%%%%%%%%%%%%%%%%%%%%%%
% %%%%%%%%%%%%%%%%%%%%%%%%%%%%%%%%%%%%%%%%%%%%%%%%%%%%%%%%%%
\pretitle{\begin{center}\bfseries\LARGE}
\posttitle{\end{center}}



{\setstretch{.8}
\maketitle
}



\section*{Introduction}

The volume of data available to investors has expanded dramatically over the past decade.
Beyond traditional firm-level and macroeconomic data, market participants increasingly rely on alternative sources such as text, web traffic, satellite imagery, and social media \citep{DataSphere2025}.
Investors now widely employ machine learning methods to extract signals from such high-dimensional data \citep{Mercer2024AI}.
A key feature of these approaches is that they enable a systematic search for predictive relationships by allowing the data itself to determine the forecasting model, without requiring prior theoretical restrictions on which variables should predict returns.
This behavior is not captured by standard learning-based asset-pricing models.
In much of the existing literature, agents rely on a pre-specified forecasting rule and update a limited set of parameters using variants of least squares \citep{Timmermann1993,hong2007simple}, or switch between a small number of competing low-dimensional models \citep{barberis1998model}.
By construction, these frameworks abstract from high-dimensional predictor spaces and from the model-selection step that is central to how modern investors learn from data.

To address this gap, we study an asset-pricing model in which returns are unpredictable under rational expectations, yet investors search for forecasting signals in a large set of potential predictors.
Our central question is whether this \emph{search for predictability} can account for two empirical features of return predictability: it is \emph{sparse}, in that only a small subset of predictors matters, and \emph{unstable}, in that predictability appears and disappears over time \citep{Chinco2019,Farmer2023}.
We model this behavior by endowing agents with a learning rule based on the least absolute shrinkage and selection operator (LASSO) \citep{Tibshirani1996}. 
In each period, agents estimate a linear forecasting model using LASSO, which sets most coefficients to zero and keeps only predictors with strong estimated effects that survive penalization. 
When LASSO selects a non-zero predictor, investors trade on the resulting perceived predictability.
Through the resulting feedback between beliefs and prices, the model generates self-sustained and short-lived windows of return predictability even though returns are unpredictable under rational expectations. 
We formally characterize when such windows arise, showing that they are more likely under higher return volatility, limited agent memory, and stronger beliefs in predictability.
We also test the model on the universe of U.S. listed stocks, combined with a large set of potential signals derived from U.S. financial and business news \citep{bybee2024business}. 
In line with our narrative, the model explains an economically meaningful share of realized returns.

In what follows, I briefly outline the model and its empirical implementation, and conclude by highlighting the main contributions and planned extensions.

\section*{Model}

\paragraph{Environment.}
We study a standard \cite{Lucas1978} asset-pricing model.
A representative risky asset pays dividend $D_t$ and its price $P_t$ satisfies the Euler equation
\begin{equation}\label{eq:euler}
P_t = \delta\,\mathbb{E}_t\!\left[\left(\frac{C_{t+1}}{C_t}\right)^{-\gamma}(P_{t+1}+D_{t+1})\right],
\end{equation}
where $C_t$ denotes consumption, preferences exhibit constant relative risk aversion with coefficient $\gamma$ and $\delta$ denotes the discount factor.
Under rational expectations, consumption and dividends follow stationary log-normal processes and returns are conditionally unpredictable.

\paragraph{Perceived law of motion.}
Following \cite{Adam2016}, we assume that agents are fully rational with respect to the dividend and consumption processes, but not with respect to returns. 
Instead, they hold the belief that returns are predictable conditional on one or more observable signals. 
We formalize this belief by endowing agents with a perceived law of motion (PLM) for log returns of the form
\begin{equation}
r_{t+1} = X_t'\beta + u_{t+1},
\end{equation}
where $X_t$ is a potentially high-dimensional predictor matrix.
Agents try to estimate the unknown coefficient vector $\beta$ using observed data.

\paragraph{Actual law of motion.}
The agents expectations of future returns $X_t'\beta$ affect actual returns trough the Euler equation.
Substituting the agents subjective beliefs into (\ref{eq:euler}), we obtain the actual law of motion (ALM) for returns: 
\begin{equation}\label{eq:alm_univariate}
    r_{t+1} = \log(\varepsilon_{t+1}) + \log(1 - \kappa e^{x_{t} \beta}) - \log(1 - \kappa e^{x_{t+1} \beta}),
\end{equation}
where $\kappa := \delta a^{- \gamma} e^{\sigma^2_u / 2}$ with $\sigma^2_u$ being the variance of $u_{t+1}$.
Note, that the resuling ALM is nonlinear, implying that the agents’ linear PLM is misspecified.

\paragraph{Equilibrium.}
When the perceived law of motion (PLM) is misspecified, learning need not converge to the rational expectations equilibrium \citep{Honkapohja2003}. 
We therefore adopt a cross-moment consistent equilibrium (CMCE), which requires internal consistency between beliefs and outcomes: the joint distribution of returns and predictors implied by the actual law of motion is such that, when agents apply their learning rule, they recover the same coefficients they use to form beliefs. 
In the CMCE, prices are determined by agents who optimize given their subjective beliefs.
These beliefs are updated using data generated by equilibrium prices, and the forecasting moments perceived by agents coincide with those implied by the actual law of motion.

\paragraph{Learning mechanism.}

Agents form expectations by estimating their PLM using the LASSO.
This regularization method allows them to consider a large set of potential predictors that can even exceed the number of available observations because LASSO shrinks small coefficients exactly to zero.
In each period, agents run a LASSO regression on a short rolling window of recent data to select a subset of predictors, estimate their coefficients, and form return expectations.
Formally, they choose the coefficient vector $\beta$ as the solution to the penalized least squares problem:
\begin{equation}
\hat{\beta}_{\text{LASSO}}  = \arg\min_{\beta}\left\{\frac{1}{2}\sum_{i=t-n}^{t-1}\big(r_{i+1}-X_i'\beta\big)^2
\;+\;\lambda \|\beta\|_1\right\},
\end{equation}
where $n$ denotes the window length and $\lambda$ the $\ell 1$-penalty. 
When regressors are orthonormal, we can derive a closed form solution that reduces to soft-thresholding of the least-squares estimate,
\begin{equation}\label{eq:lasso_from_ols}
\hat{\beta}_{\text{LASSO}} = S(\beta_{\text{OLS}}, \lambda) = \textrm{sign}(\beta_{\text{OLS}}) \cdot \max\left(0, |\beta_{\text{OLS}}| - \lambda\right),
\end{equation}
As a result, a predictor enters the forecasting model only if the coefficient agents would estimate using ordinary least squares (OLS) exceeds the penalty $\lambda$ in absolute value.

\paragraph{Main result.}

Under LASSO learning, asset returns can become predictable even though they are unpredictable under rational expectations.
This predictability arises endogenously from the feedback between investors’ beliefs and equilibrium prices.
We characterize the conditions under which such episodes occur by deriving the probability that a given predictor is selected by the LASSO.
To obtain a tractable characterization of this mechanism, we assume that agents update their beliefs using constant-gain learning with gain parameter $g$.\footnote{In the empirical application, we estimate the model using a rolling window that places equal weight on a fixed number of recent observations. Since this approach is analytically intractable, we approximate it in the theory with constant-gain learning. The two schemes can be related by matching the cumulative weight assigned to recent data: letting $p$ denote the proportion of total weight placed on the $\ell$ most recent observations under constant-gain learning with decay parameter $g$, the implied window length satisfies $\ell \approx \log(1-p)/\log(1-g)$.}
For $\kappa < 1$ and small values of $g$, the probability that the LASSO estimate is nonzero equals the probability that the corresponding ordinary least squares estimate exceeds the penalty parameter $\lambda$:
\begin{equation}
\mathbb{P}\big(|\beta_{\text{OLS}}| > \lambda\big)
= 1 - \left[
\Phi\left(\frac{\lambda}{\sigma_{\text{OLS}}}\right)
- \Phi\left(\frac{-\lambda}{\sigma_{\text{OLS}}}\right)
\right],
\end{equation}
where
\begin{equation}
\sigma_{\text{OLS}}
= \left(
g s_d^2
\left(
2 \sigma_x^2
\left(1 + \frac{\kappa}{1 - \kappa}\right)
\right)^{-1}
\right)^{1/2}.
\end{equation}


\paragraph{Estimation}

To estimate the model, we use two data sources.
First, we collect daily returns for all U.S. listed stocks from CRSP over the period 1984-2017.
Second, we use the daily attention of news to 180 financial and business topics constructed by \cite{bybee2024business} from a large corpus of U.S. news articles.
We implement the model in a two-stage empirical design.
In the first stage, we estimate a rolling LASSO forecasting regression using the news attention series as candidate predictors.
In the second stage, we use the resulting return forecasts to estimate the ALM.
If our postulated mechanism is at work in the data, the return forecasts from the first stage should explain a significant share of realized returns in the second stage.
We show that this is indeed the case: the return forecasts explain an economically meaningful share of realized returns.


\section*{Contribution}

This paper makes three main contributions to the literature in asset pricing.
First, in learning-based asset-pricing models, investors typically update beliefs within a small, pre-specified forecasting rule \citep{Marcet1989,Timmermann1993,barberis1998model,hong2007simple}.
We instead study a setting in which investors have no prior idea of the true return process but instead select a new model each time they update beliefs. 
This novel learning mechanism is a closer approximation of how investors today utilize data.
Second, building on the internal rationality approach of \cite{Adam2016} in the Lucas-tree tradition \citep{Lucas1978}, we allow beliefs and prices to interact: when investors come to believe a signal predicts returns, they trade on it, and this trading feeds back into prices.
This feedback loop between statistical learning and prices is new to the literature.
Third, the model offers a parsimonious explanation for the sparse and episodic nature of return predictability documented in recent empirical work \citep{Chinco2019,Farmer2023}.
In our framework, predictability arises intermittently from selection events and fades as signals fail to survive re-estimation and regularization. 
This mechanism delivers sharp comparative statics, linking the frequency of predictability episodes to market conditions and the characteristics of agents search behavior.

\section*{Outlook}

We plan to extend our model in several directions.
First, we plan to introduce heterogeneous agents who differ in their strength of regularization.
For example, optimistic investors may penalize weak signals less aggressively, while pessimistic investors apply stronger penalties, generating heterogeneity in perceived predictability and trading behavior.
Second, we will broaden the set of candidate predictors to include alternative classes of signals, such as firm fundamentals, technical indicators, and additional forms of textual and non-textual information.
Third, we will study the role of memory more systematically by allowing agents to differ in the length of the estimation window.
These extensions will allow us to assess how heterogeneity in beliefs, information, and memory shapes the dynamics of predictability and prices.

\newpage

\bibliography{references}

\end{document}